\documentclass[11pt,reqno]{amsart}
\usepackage[T1]{fontenc}
\usepackage{lmodern,amsmath,amssymb,mathtools,microtype,booktabs,array}
\usepackage[margin=1in]{geometry}
\usepackage[hidelinks]{hyperref}
\hypersetup{pdftitle={General Theta Foundations I: Acquired Geometry and Causal Resource Transfer},pdfauthor={Qian Qi}}
\ifdefined\pdfinfoomitdate\pdfinfoomitdate=1\fi
\ifdefined\pdftrailerid\pdftrailerid{}\fi
\ifdefined\pdfsuppressptexinfo\pdfsuppressptexinfo=15\fi
\numberwithin{equation}{section}
\newtheorem{theorem}{Theorem}[section]
\newtheorem{proposition}[theorem]{Proposition}
\newtheorem{lemma}[theorem]{Lemma}
\newtheorem{corollary}[theorem]{Corollary}
\theoremstyle{definition}
\newtheorem{definition}[theorem]{Definition}
\newtheorem{example}[theorem]{Example}
\theoremstyle{remark}
\newtheorem{remark}[theorem]{Remark}
\newcommand{\E}{\mathbb E}
\newcommand{\Pp}{\mathbb P}
\newcommand{\R}{\mathbb R}
\newcommand{\N}{\mathbb N}
\newcommand{\one}{\mathbf1}
\newcommand{\norm}[1]{\lVert#1\rVert}
\newcommand{\ip}[2]{\langle#1,#2\rangle}
\newcommand{\tr}{\operatorname{tr}}
\newcommand{\law}{\operatorname{Law}}
\newcommand{\conv}{\operatorname{conv}}
\newcommand{\supp}{\operatorname{supp}}
\newcommand{\TV}{\mathrm{TV}}
\newcommand{\cp}{\mathrm{cp}}
\newcommand{\aut}{\mathrm{aut}}
\newcommand{\calR}{\mathcal R}
\newcommand{\calF}{\mathcal F}
\newcommand{\calA}{\mathcal A}
\newcommand{\eps}{\varepsilon}
\newcommand{\HS}{\mathrm{HS}}
\emergencystretch=2em
\title[Acquired geometry and causal resource transfer]{General Theta Foundations I:\\Acquired Geometry and Causal Resource Transfer}
\author{Qian Qi}
\date{October 2026}
\subjclass[2020]{Primary 60G40, 93E11; Secondary 28A80, 62C20, 81P15, 94A29}
\keywords{prepared experiment, predictive quotient, adaptive acquisition, finite-state observer, stopping, actual occupation, singular quantization}
\begin{document}

\begin{abstract}
We prove a resource--risk theorem for prepared finite-action causal experiments admitting a history-uniform calibration certificate; the current report need not recover the posterior. A kernel-level calibration certificate couples an observable prediction proxy to a nonnegative residual-risk mark. The actual law of these two quantities defines a marked occupation profile. It bounds every history-dependent acquisition and stopping policy from below and is implemented, up to the stated calibration constants, by one live state and the allotted output labels. If location mass has exponent $d$ and low-uncertainty mass has exponent $\alpha$, the matching law is the larger of $N^{-2/\alpha}$ and $(Nm)^{-2/(d+\alpha)}$; the crossover is derived from actual mass, not an ambient dimension. Correlated hidden-state sensors and noncommuting quantum instruments verify the same theorem. A Gaussian consequence gives a matching acquisition rate for the unrestricted-history optimum and a three-state observer, despite nonregeneration. A marked revealing component produces an exact zero-risk transition; for two-state Gaussian observers we determine the complete curve and its different critical budget. Protected calibration, causal simulation, numerical precision and all observer resources are separately charged.
\end{abstract}
\maketitle
\section{Introduction}
An observation can be favorable in two different ways: it can lie near a representable prediction, and it can leave little uncertainty about the protected future test. These are not the same property. In an informative hidden-state experiment, the last report usually does not determine the posterior. Compressing a convenient function of that report while ignoring the remaining conditional variance can therefore give an arbitrarily optimistic resource law.

We study the experiment-level passage
\[
\begin{aligned}
&\text{preparation and positive kernels}\\
&\quad\longrightarrow\text{executable future tests and predictive quotient}\\
&\quad\longrightarrow\text{actual marked acquisition law}\\
&\quad\longrightarrow\text{causal implementation and protected-task risk}.
\end{aligned}
\]
The principal theorem uses a calibrated report proxy $f$ together with a residual-risk mark $w\ge0$. Its assumptions do not assert that $f$ is a sufficient statistic. Instead, at every legal history, the conditional variance above a common irreducible baseline is comparable to $w$, and the squared bias of $f$ is bounded by $w$. These inequalities have an unnormalized one-call formulation in the raw instrument. Consequently they can be checked on an invariant family of prepared physical laws, before any observer or budget is selected.

The resulting geometry is a probability measure on pairs $(f,\sqrt w)$, with distortion $w+\min_j\|f-r_j\|^2$. This is an actual one-call law of a specified prepared reference experiment. Every adaptive stopping procedure induces a terminal submeasure bounded by its paid acquisition occupation. Conversely, a common parameter-independent program uses fresh action and acceptance coins and only one live label to realize the corresponding upper bound. Histories continue to alter the hidden state and the conditional audit mean; no posterior reconstruction from a last report is presumed. Theorem~\ref{thm:central} supplies both directions and an explicit mass--cost dual.

When location and low-uncertainty masses have exponents $d$ and $\alpha$, respectively, the profile has two distinct constraints. Concentrating all terminal mass on small uncertainty costs $N^{-2/\alpha}$. Simultaneously concentrating it near one of $m$ output values costs $(Nm)^{-2/(d+\alpha)}$. Theorem~\ref{thm:geometry} proves that the larger term is both necessary and attainable. For an autonomous observer, $m=M-1$ counts every output label after paying for the live controller. The resulting saturation threshold $m\asymp N^{d/\alpha}$ is a statement about the acquired task geometry, not about the dimension of the physical carrier or the number of stored real coordinates.

Two informative realizations verify this same general theorem. The first carries a genuinely evolving hidden Markov residual through controlled observations. The second retains quantum coherence under noncommuting weak instruments. Both include a reported location and reliability, while the unrevealed task component remains history dependent. Singular regular supports, including Cantor products, are permitted without a smooth manifold assumption.

A separate consequence is a natural Gaussian benchmark. For a hidden sign observed through $N(\pm\mu,1)$ likelihoods, an all-history converse and a three-state construction have the same acquisition rate
\[
 \log R(N)=-2\mu\sqrt{2\log N}+O(1).
\]
The lower bound allows arbitrary stopping and arbitrary terminal readouts; it does not compare only one controller. A revealing branch creates a non-power-law critical transition. In the symmetric classical chain we additionally determine the exact risk of every optimal two-state observer, including randomized rules, and its larger zero-risk budget. These are consequences of the general risk-envelope mechanism and actual occupation balance, not substitutes for the general theorem.

\subsection*{Relation to earlier work}
Conditional projection, randomized selection by conditional risk, and finite-state likelihood-ratio testing are established tools. Selective classification optimizes uncertainty versus coverage~\cite{Franc}; trimmed quantization optimizes a bounded density and a codebook~\cite{CGM}; finite-memory hypothesis testing has a classical theory~\cite{HC}. We do not claim these mechanisms as new. The present comparison concerns a changing protected state, history-dependent report laws, a paid stopping rule, and a complete autonomous state count. Its specific new conclusions are the calibrated marked occupation transfer, the joint reliability--resolution crossover, and the stated informative and singular acquisition laws.

Predictive equivalence~\cite{SC}, value-directed belief compression~\cite{PB}, and Bayes-risk/Bregman geometry~\cite{Varshney} concern related but different information objects. The current finite-memory belief approximation result~\cite{Kim} explicitly compares costs under the same closed-loop execution; our lower bounds quantify over every legal history-dependent stopping policy. Gaussian filtering and weak quantum measurement supply the raw experimental context~\cite{Rabiner,Jacobs}. The retained foundations articles contain the exact common-program occupation criterion, the likelihood-aware three-defect identity, and the report-recoverable acquisition theorem~\cite{GTF49,GTF32,GTF51}. We preserve those results with their original hypotheses. Here the risk mark removes the report-recoverability restriction for the declared calibrated class; it does not claim a necessary and sufficient classification of all causal experiments.

\section{Prepared experiments and calibration}
\label{sec:experiment}
All spaces are standard Borel. A preparation and positive normalized kernels
\[
 K_a^\theta(dx',dy\mid x),\qquad a\in\mathcal A,
\]
are given; $\mathcal A$ is finite and nonempty. The fixed model parameter $\theta$ is unknown to the observer. All actions, reports, failed acquisitions and physical costs are part of the experiment. A bounded audit $Z\in\R^k$ is performed after the observer commits to its response, without a subsequent uncontrolled state transition. Its distribution is specified by the stopped physical state. Let $H\subset\R^k$ be a known compact convex set containing every audit value and every proxy used below. Predictions outside $H$ may be projected into $H$.

Write $\mathcal H_{t-1}$ for the complete visible history and all previous private randomizations, before action $A_t$ is selected. At a surviving call, define $m_{t,a}(y)$ and $v_{t,a}(y)$ to be the conditional audit mean and total conditional variance given $(\mathcal H_{t-1},A_t=a,Y_t=y)$. These are regular conditional quantities; no positive mass is assigned to a continuous singleton. Fresh routing coins are independent of the audit conditional on that information.

\begin{definition}[Raw calibrated proxy certificate]\label{def:certificate}
Known Borel functions $f_a:Y_a\to H$, $w_a:Y_a\to[0,W]$, actual reference one-call laws $\nu_a$, constants $B_0\ge0$, $c_v>0$, $C_v,c_b\ge0$, and $0<\ell_a\le1\le u_a<\infty$ satisfy, for every model, every time and every legal surviving history,
\begin{align}
 \ell_a\nu_a&\le\law_\theta(Y_t\mid\mathcal H_{t-1},A_t=a)\le u_a\nu_a,
 \label{eq:report-envelope}\\
 c_vw_a(y)&\le v_{t,a}(y)-B_0\le C_vw_a(y),\qquad
 \|m_{t,a}(y)-f_a(y)\|^2\le c_bw_a(y).
 \label{eq:calibration}
\end{align}
The inequalities hold on the common reference-almost-everywhere domain appropriate to each conditional law. Each $\nu_a$ must be obtained from a specified preparation and action in an admitted reference member. The current proxy is not asserted to be the full predictive quotient.
\end{definition}

Here is a sufficient one-call verification that avoids postulating a property of a chosen observer. Let $\mathcal B$ be a family of physical probability laws containing the preparations and closed under every normalized report instrument. For $b\in\mathcal B$, let $\lambda_a(y;b)$, $U_a(y;b)$ and $V_a(y;b)$ be the unnormalized report density, audit first moment and audit squared-norm moment, relative to $\nu_a$. It suffices to check, uniformly in $b$ and $\theta$,
\begin{align}
 &\ell_a\le\lambda_a\le u_a,\label{eq:raw-certificate}\\
 &c_vw_a\lambda_a^2\le V_a\lambda_a-\|U_a\|^2-B_0\lambda_a^2
       \le C_vw_a\lambda_a^2,\notag\\
 &\|U_a-\lambda_af_a\|^2\le c_bw_a\lambda_a^2.\notag
\end{align}
Indeed division by $\lambda_a^2$ gives \eqref{eq:calibration}; induction through the raw kernels gives every history quantifier. No Fisher matrix, covariance geometry, smooth state parameterization or contraction rate is an input. The variance in \eqref{eq:raw-certificate} is derived from the specified audit.

A checkpoint may use the entire history to choose actions and a randomized stopping time $1\le\tau<\infty$, then encode it once into at most $m$ fixed output values. An autonomous observer has disjoint live and terminal label sets $S,J$, a fixed initial live label, an action row $p(a\mid s)$, a Borel routing row $\gamma(s'\text{ or }j\mid s,a,y)$, and a fixed output $r_j$. All these objects are parameter independent. Its budget is $|S|+|J|\le M$. Retained seeds, free clocks, old reports and exact posterior registers are forbidden; fresh independent randomization is allowed. Preparation, audit, program description and temporary computation are counted separately below.

For known positive action costs $c_a$, the common budget is
\[
 \sup_\theta\E_\theta\sum_{t=1}^{\tau}c_{A_t}\le N.
\]
The resulting worst-model risk infima are $R_{\cp}^\Theta(N,m;c)$ and $R_{\aut}^\Theta(N,M;c)$. With unit costs this is a mean raw-call budget, not a deterministic deadline. The live-label price is therefore one, not the fixed-horizon phase deduction from a different interface. Every lower bound below permits arbitrary history-dependent actions; every upper is one common finite-state program.

\section{The marked acquired-geometry transfer theorem}
\label{sec:central}
For $C=(r_1,\ldots,r_m)\in H^m$, define
\[
 d_{a,C}(y)=w_a(y)+\min_j\|f_a(y)-r_j\|^2.
\]
The marked occupation profile is
\begin{align}
 \Xi_m(N;c)=\min\ &\sum_a\int d_{a,C}(y)h_a(y)\nu_a(dy),\label{eq:marked-profile}\\
 &C\in H^m,\quad x_a\ge0,\quad0\le h_a\le x_a\quad\nu_a\text{-a.e.},\notag\\
 &\sum_a\int h_a\,d\nu_a=1,\qquad \sum_ac_ax_a\le N.\notag
\end{align}
Infeasible values are $+\infty$. The unnormalized measure $h_a\nu_a$ is selected terminal mass, and $x_a$ is an expected paid action count in the reference allocation problem. It is not automatically the actual terminal measure of a correlated implementation.

Put
\[
 k_- =\left[\max\left\{2,\frac{1+2c_b}{c_v}\right\}\right]^{-1},
 \qquad k_+=\max\{2,C_v+2c_b\},\qquad
 K_0=\max_a\frac{u_a}{\ell_a}.
\]
Division of a cost vector by $u$ or $\ell$ is componentwise.

\begin{theorem}[Causal transfer through marked acquired mass]\label{thm:central}
Under Definition~\ref{def:certificate}, for every budget and every $m\ge1$, $M\ge2$,
\begin{align}
 k_-\Xi_m(N;c/u)
 &\le R_\cp^\Theta(N,m;c)-B_0
 \le k_+K_0\Xi_m(N;c/\ell),\label{eq:cp-transfer}\\
 k_-\Xi_{M-1}(N;c/u)
 &\le R_\aut^\Theta(N,M;c)-B_0
 \le k_+K_0\Xi_{M-1}(N;c/\ell).\label{eq:aut-transfer}
\end{align}
An upper with an infeasible profile is vacuous. Each finite upper is implemented with one live label and the prescribed output labels, at the original budget $N$, uniformly over the unknown model. No stability of a stored approximate posterior is required.

The profile minimum exists whenever $N\ge\min_a c_a$. For fixed $C$ its value is
\begin{equation}\label{eq:dual}
 \sup_{z\in\R,\lambda\ge0}\left\{z-\lambda N:
       \int(z-d_{a,C})_+\,d\nu_a\le\lambda c_a\quad\text{for every }a\right\}.
\end{equation}
The supremum is attained for $N>\min c_a$; at the endpoint only equality is asserted. The minimizing allocation may be chosen to use at most two actions. The codebook is optimized outside the displayed dual, not exchanged with its supremum.

If $w_a=0$, the conditional means equal $f_a$, the variance is $B_0$, and the conditional report laws equal $\nu_a$, then the stronger exact equalities hold:
\[
 R_\cp(N,m;c)-B_0=\Xi_m(N;c),\qquad
 R_\aut(N,M;c)-B_0=\Xi_{M-1}(N;c).
\]
Thus the report-recoverable theorem is a special case, not a separate replacement for the present result.
\end{theorem}

The pushforward of $\nu_a$ by $y\mapsto(f_a(y),\sqrt{w_a(y)})$ is the \emph{marked acquired law}. Its support and small-ball masses are the geometric objects used below. This construction is task dependent: it makes no assertion that a finite task proxy is a sufficient state for every executable future experiment. It also makes no claim that a variance proxy alone implies compressibility; both inequalities in \eqref{eq:calibration} are necessary for the stated comparison mechanism.

The two cost scalings in \eqref{eq:cp-transfer}--\eqref{eq:aut-transfer} are deliberate. Uniform report domination does not by itself imply that $\Xi(cN)$ and $\Xi(N)$ have a bounded ratio at a critical atom. Matching at the same asymptotic scale is proved under explicit geometric hypotheses in the next sections; the singular critical point is analyzed separately without rescaling away its actual mass.

\section{Proof of the transfer theorem}
\begin{lemma}[Conditional task comparison]\label{lem:comparison}
For every legal pre-audit history, report and prediction $r\in H$,
\[
 k_-\{w_a(y)+\|f_a(y)-r\|^2\}
 \le \E[\|Z-r\|^2\mid\mathcal H_{t-1},a,y]-B_0
 \le k_+\{w_a(y)+\|f_a(y)-r\|^2\}.
\]
The same inequalities remain valid after an independent routing coin chooses $r$.
\end{lemma}
\begin{proof}
Write $m=m_{t,a}(y)$, $v=v_{t,a}(y)$ and $w=w_a(y)$. Conditional projection gives the middle expression as $v-B_0+\|m-r\|^2$. The upper follows from
$\|m-r\|^2\le2\|f_a-r\|^2+2c_bw$.
For the lower,
\[
 w+\|f_a-r\|^2\le(1+2c_b)w+2\|m-r\|^2
 \le k_-^{-1}\{v-B_0+\|m-r\|^2\}.
\]
The coin has no audit information, so conditioning on its selected value does not change $m$ or $v$.
\end{proof}

\begin{lemma}[Actual stopped occupation]\label{lem:occupation}
For any legal history-dependent procedure of finite expected cost, put
\[
 x_a=\E\sum_{t\le\tau}\one_{\{A_t=a\}},\qquad
 \zeta_a(D)=\Pp(A_\tau=a,Y_\tau\in D).
\]
Then $\sum_a\zeta_a(Y_a)=1$, $\sum_ac_ax_a\le N$, and
$\zeta_a\le u_ax_a\nu_a$. These are assertions about its actual execution law.
\end{lemma}
\begin{proof}
Positive minimum cost implies $\E\tau<\infty$ and hence almost-sure termination. For a Borel set $D$, the expected number of action-$a$ reports in $D$ is
\[
 \sum_t\E\left[\one_{\{\tau\ge t,A_t=a\}}
       \Pp(Y_t\in D\mid\mathcal H_{t-1},A_t=a)\right]
 \le u_ax_a\nu_a(D).
\]
The terminal count is no larger than this acquired count. Predictability is used before the report; conditioning the law after rejection back to a reference law would not justify this calculation.
\end{proof}

\begin{proof}[Lower bounds]
Apply Lemma~\ref{lem:comparison} to the actually committed output, then replace its squared distance by the minimum over the fixed output dictionary. With $h_a=d\zeta_a/d\nu_a$ and $x'_a=u_ax_a$, Lemma~\ref{lem:occupation} yields a feasible point of $\Xi_m(N;c/u)$. This proves the checkpoint lower for each model and then for the common worst-model infimum. An autonomous program has at least one live label, so at most $M-1$ terminal values. Its entire action and stopping policy is already included in the preceding lower bound.
\end{proof}

\begin{proof}[Common finite-state upper]
Take a minimizing point $(C,x,h)$ of $\Xi_m(N;c/\ell)$ and put $T=\sum_a x_a/\ell_a\ge1$. At one live label, choose action $a$ with probability $x_a/(\ell_aT)$. For $x_a>0$, accept report $y$ with probability $h_a(y)/x_a$, route to a nearest proxy center in $C$, and otherwise continue at the same live label. Actions with $x_a=0$ are not selected. A fixed Borel tie rule is used.

At every surviving history, the conditional probability of stopping is at least
$T^{-1}\sum_a\int h_a\,d\nu_a=T^{-1}$. Therefore $\E_\theta\tau\le T$, uniformly in $\theta$. Conditional expected acquisition cost per call is the fixed number $T^{-1}\sum_ac_ax_a/\ell_a$, so total expected cost is at most $N$. The conditional contribution to protected excess risk on one call is at most
\[
 \frac{k_+}{T}\sum_a\frac{u_a}{\ell_a}\int d_{a,C}h_a\,d\nu_a
 \le \frac{k_+K_0}{T}\Xi_m(N;c/\ell).
\]
Summing over actual survival probabilities proves the upper. It uses no persistent real vector, posterior or random seed. The parameter-independent program has one live label and $m$ terminal labels; its actual selected law need not equal $h_a\nu_a$.
\end{proof}

\begin{proof}[Existence, dual and exact special case]
The counts belong to $0\le x_a\le N/c_a$. On this box the functions $h_a$ lie in weak-star compact subsets of $L^\infty(\nu_a)$. The inequalities $0\le h_a\le x_a$ and the mass constraint are weak-star closed, as can be tested against nonnegative $L^1$ functions. The centers belong to compact $H^m$, and their squared distance functions converge uniformly when the centers converge. Since $w_a$ is bounded, the objective is jointly continuous in these topologies. This proves existence and the feasibility criterion.

For fixed centers, minimizing the mass--cost Lagrangian over $0\le h_a\le x_a$ gives
\[
 z-\lambda N+\sum_ax_a\left\{\lambda c_a-\int(z-d_{a,C})_+d\nu_a\right\}.
\]
Minimizing in $x_a\ge0$ yields \eqref{eq:dual}. To prove equality, consider the convex cone of achievable triples (terminal mass, cost, marked distortion), with nonnegative slack in the last two coordinates. Its bounded-cost slices are closed by the compactness just proved. At mass one and $N>\min c_a$, the projected point is interior to the feasible mass--cost cone. A supporting plane to the finite convex distortion value has a nonpositive cost coefficient and can be normalized to distortion coefficient one. Testing it on the one-action rays gives precisely \eqref{eq:dual}, with equality at the optimum. At the endpoint, take budgets decreasing to $\min c_a$; compactness gives convergence of the primal values, and evaluation of the interior dual optimizers at the endpoint gives the supremum formula. No boundedness of endpoint multipliers is inferred.

For the two-action assertion, at an optimizer discard actions with no terminal mass. Freeze each remaining action's conditional terminal distribution, per-unit terminal cost and per-unit terminal distortion. Reoptimize only its weights subject to total mass one and one cost inequality. An extreme minimizer has at most two positive weights: three positive weights admit a nonzero two-sided perturbation preserving total mass and cost. Rescaling the frozen counts and densities preserves their pointwise domination. This standard finite-dimensional argument makes no claim about a uniquely optimal action mixture.

Finally, in the zero-mark regenerative case the conditional loss is exactly $B_0+\|f_a-r\|^2$. The lower occupation inequality uses $u_a=1$. The construction chooses $a$ with probability $x_a/\sum x_a$ and accepts with $h_a/x_a$; its stopping mass is exactly $1/\sum x_a$ at every call. Its actual terminal submeasures are exactly $h_a\nu_a$, and its expected action counts are exactly $x_a$. The exact equalities follow. This last mechanism is the inherited report-recoverable result~\cite{GTF51}; the conditional task comparison is what permits the nonzero-mark extension.
\end{proof}

\section{Joint location and reliability geometry}\label{sec:geometry}
We first use one unit-cost action. Let $\sigma$ be the actual pushforward of $\nu$ under $(f,\sqrt w)$. Suppose $f\in K$, $0\le\sqrt w\le1$, where $K\subset H$ is compact, and suppose there are $d,\alpha>0$, $r_0>0$ and $0<c_0\le C_0<\infty$ such that
\begin{equation}\label{eq:rectangle-mass}
 c_0r^ds^\alpha\le
 \sigma\bigl(B(x,r)\times[0,s]\bigr)
 \le C_0r^ds^\alpha
 \quad(x\in K,
\ 0<r\le r_0,
\ 0<s\le1).
\end{equation}
These are masses of acquired rectangles, not volume assigned to the formal physical state space. Independence of $f$ and $w$ is sufficient but is not required. Finitely many actions may satisfy these same inequalities with uniform constants; the result below remains valid because one action supplies the upper and summed occupations supply the lower.

\begin{theorem}[Acquisition--resolution crossover]\label{thm:geometry}
Under \eqref{eq:rectangle-mass}, uniformly for every $N\ge1$ and $m\ge1$,
\begin{equation}\label{eq:profile-crossover}
 \Xi_m(N)\asymp
 F_{d,\alpha}(N,m):=
 \max\{N^{-2/\alpha},(Nm)^{-2/(d+\alpha)}\}.
\end{equation}
Consequently, under the calibration and report-envelope conditions with unit costs,
\begin{align}
 R_\cp^\Theta(N,m)-B_0&\asymp F_{d,\alpha}(N,m),\label{eq:crossover-cp}\\
 R_\aut^\Theta(N,M)-B_0&\asymp F_{d,\alpha}(N,M-1),\qquad M\ge2.
 \label{eq:crossover-aut}
\end{align}
Constants may depend on the geometric and calibration constants and on the finite action family, but not on $N,m,M$, the chosen policy or the model. The bounds are simultaneously information-theoretic and implementable in the declared finite-label model.
\end{theorem}
\begin{proof}
The marginal $f_\#\nu$ is $d$-regular by setting $s=1$. Cover $K$ by finitely many radius-$r_0$ balls and use the upper bound to obtain
$\sigma(K\times[0,s])\le C s^\alpha$. Any feasible terminal density is at most $N$. Choosing $s=cN^{-1/\alpha}$ with $c$ small leaves at least half its terminal mass at $w>s^2$. This gives the first lower term.

For any $m$ centers, the sublevel $w+\min_j\|f-r_j\|^2\le t^2$ lies in the union of $m$ rectangles of radius $t$ and height $t$. Centers outside $K$ can be moved to points of a nonempty intersecting ball, doubling the radius. Its reference mass is at most $Cmt^{d+\alpha}$, and terminal mass is at most $CNmt^{d+\alpha}$. With $t=c(Nm)^{-1/(d+\alpha)}$, at least half the terminal mass remains outside. Adjusting $c$ ensures every radius is within $r_0$; bounded small budgets only change the fixed constant. This proves the second lower term.

For the upper, $d$-regularity gives both global covering and packing. A maximal $r$-separated set is an $r$-cover, and its disjoint $r/3$ balls give cardinality at most $Cr^{-d}$. Conversely, a greedy packing contains $k$ points at pairwise distances at least $c k^{-1/d}$, since otherwise the upper mass bound on fewer than $k$ such balls would fail to cover probability one.

For sufficiently large $N$, choose
$k=\min\{m,\lfloor c_1N^{d/\alpha}\rfloor\}$, replacing a zero floor by one. Choose the small constant $c_1$ and then
$t=C_1(Nk)^{-1/(d+\alpha)}$ so that the $k$ radius-$t$ balls around the packing points are disjoint and the union of their marked rectangles of height $t$ has mass at least $1/N$. The two requirements are compatible because $k\le c_1N^{d/\alpha}$. Restrict to this union with a constant density giving terminal mass one; that density is at most $N$. Its distortion is at most $2t^2$. Since $k\asymp\min(m,N^{d/\alpha})$, this is the right side of \eqref{eq:profile-crossover}. Bounded $N$ is handled by the whole law and one fixed center; its distortion is bounded and $F_{d,\alpha}$ is bounded below there.

For several uniformly certified actions, the first lower calculation uses the total acquired count, and the second uses the sum of the action-indexed masses, again at most $N$. One chosen action gives the upper. Finally $F(cN,m)\asymp F(N,m)$ for every fixed $c>0$ where both budgets are admissible. Apply Theorem~\ref{thm:central}; when $\ell N<1$, a one-call fixed decoder and bounded loss cover the remaining bounded interval. This proves \eqref{eq:crossover-cp}--\eqref{eq:crossover-aut}.
\end{proof}

\begin{remark}[Why both terms are necessary]\label{rem:two-terms}
The estimate $(Nm)^{-2/(d+\alpha)}$ alone would tend to zero at fixed $N$ as $m\to\infty$. That is impossible here: even a decoder with the complete history must pay to acquire low conditional variance. Conversely, the term $N^{-2/\alpha}$ alone misses the finite dictionary cost when $m\ll N^{d/\alpha}$. No optimal allocation claim is hidden in a fixed-dimensional constant.
\end{remark}

\subsection*{Singular supports}
Let $K$ be a finite product of intervals and middle-third Cantor sets, equipped with the product of uniform and fair Cantor measures. It is $d$-regular, where each interval contributes one and each Cantor factor contributes $\log2/\log3$. For a Cantor factor, choose the cylinder scale $3^{-(j+1)}<r\le3^{-j}$. A ball intersects a uniformly bounded number of level-$j$ cylinders, each of mass $2^{-j}$, giving the upper bound. A deeper cylinder through its center, of depth differing by a fixed number, lies in the ball and gives the lower bound. Product balls are sandwiched between fixed rescalings of product rectangles, proving the assertion for products. With an independent reliability variable $T$ satisfying $\Pp(T\le s)=s^\alpha$, the actual marked law satisfies \eqref{eq:rectangle-mass}. Thus noninteger acquired dimensions and singular preparation measures enter the same proved theorem, without ambient whitening or a tangent-rank assumption.

\section{Two informative raw-kernel realizations}\label{sec:realizations}
Both constructions start with a compact $d$-regular probability $\nu_K$ and a reliability law $\rho_\alpha(dt)=\alpha t^{\alpha-1}dt$ on $(0,1]$. The measure may have the singular location support just described. Their actual prepared reference reports have law
\[
 \nu(dy,dt,dv)=\nu_K(dy)\rho_\alpha(dt)\tfrac12(\delta_{-1}+\delta_1)(dv).
\]
The common proxy and risk mark are $f(y,t,v)=y$ and $w(y,t,v)=t^2$. The protected task is vector prediction, not parameter estimation.

\begin{theorem}[A controlled hidden-state sensor]\label{thm:classical}
Let a hidden sign $E$ evolve, before each report, by one of finitely many controlled transition matrices. Every entry lies in $[q,1-q]$, where $0<q\le1/2$; the matrices may belong to an unknown fixed family. Independently draw $(Y,T)$ with law $\nu_K\otimes\rho_\alpha$, and report $V\in\{-1,1\}$ with
\[
 \Pp(V=v\mid E=e)=\frac{1+\eta ve}{2},\qquad0<\eta<1.
\]
The current physical state includes $(E,Y,T)$. After stopping, audit
$Z=Y+TEe_1$, where $e_1$ is a fixed unit vector. No transition occurs before this audit. The observer sees $(Y,T,V)$, not $E$.

Theorem~\ref{thm:central} applies uniformly to the whole transition family, and the matching law \eqref{eq:crossover-aut} holds. A reference member with a symmetric transition and initial fair sign gives the stated actual first-call law. Rejected calls do not reset the sign, and the conditional audit mean is generally not a function of the last report alone.
\end{theorem}
\begin{proof}
Set $\gamma=(1-\eta)/2$ and $a=q\gamma/(1-\gamma)>0$. Before each report, the conditional probability of either sign is at least $q$, whatever the surviving history and chosen transition. Bayes' rule after $V$ gives probabilities in $[a,1-a]$. The report likelihood relative to the fair $V$ reference is between $2\gamma$ and $2(1-\gamma)$, and $Y,T$ retain their displayed conditional law. Therefore \eqref{eq:report-envelope} holds at every history with $\ell=2\gamma$, $u=2(1-\gamma)$.

Writing $b=\E[E\mid\mathcal H_{t-1},Y,T,V]$, the audit mean is $Y+Tb e_1$ and its variance is $T^2(1-b^2)$. Hence
\[
 4a(1-a)T^2\le v\le T^2,\qquad \|m-Y\|^2\le T^2.
\]
These prove the raw calibration certificate with $B_0=0$, $c_v=4a(1-a)$, $C_v=c_b=1$. They can equivalently be checked by multiplying through the positive report likelihood as in \eqref{eq:raw-certificate}. The marked law is the actual law of $(Y,T)$, so \eqref{eq:rectangle-mass} follows from $d$-regularity and the exact mass $s^\alpha$.

For completeness, the predictive carrier is obtained from the raw two-coordinate positive instrument, with the reported $Y,T$ included for a possible immediate audit. Future instruments act linearly on the unnormalized two sign masses; audit probabilities and their moments are executable functions of that state. Distinct sign beliefs are separated by a future $V$ report when the chosen transition is invertible; when it is not, the redundant direction is quotiented out rather than declared identifiable. An invertible transition exists in the nontrivial reference subclass $q<1/2$. The last reported $Y,T$ provide a Borel realization of the immediate-task coordinates; the full predictive quotient is not the proxy $Y$.

Report total variation is continuous, indeed Lipschitz, in the incoming sign belief. Bayesian updates are ratios with denominators at least $\gamma$, so are uniformly continuous on the compact belief interval. Positive transitions supply an invariant interval before each report. The risk proof uses this invariant calibration, not a stored posterior recursion or an uncharged numerical register. It remains valid for any causal transition selection respecting the same raw entry bounds.
\end{proof}

\begin{theorem}[Noncommuting weak quantum sensors]\label{thm:quantum}
Prepare a qubit in $I/2$. Let
\[
 D_q(\rho)=(1-2q)\rho+q\,\tr(\rho)I,
\qquad0<q\le1/2.
\]
For action $a$ choose a known unit vector $n_a\in\R^3$, a unitary $U_a$, and $0<\eta<1$. Write $\sigma=(\sigma_x,\sigma_y,\sigma_z)$ for the Pauli matrices and define
\[
 K_{a,v}=\left[\frac{I+\eta v\,n_a\cdot\sigma}{2}\right]^{1/2},
 \qquad
 \mathcal I_{a,v}(\rho)=K_{a,v}U_aD_q(\rho)U_a^*K_{a,v}.
\]
Independently report $(Y,T)$ with the same law as above. Retain the normalized post-instrument physical state. After commitment, measure $\sigma_z$ on that state to obtain $B\in\{-1,1\}$ and audit $Z=Y+TB e_1$.

These raw instruments verify the same Theorem~\ref{thm:central} and the same crossover \eqref{eq:crossover-aut}. The directions $n_a$ may be nonparallel and the instruments need not commute. Uniform bounds permit unknown $q$, directions or channels in a fixed admissible family; the implementing proxy program does not read the true member.
\end{theorem}
\begin{proof}
Each operation is completely positive and $\sum_vK_{a,v}^2=I$, so the instrument is normalized. On a density input, $D_q(\rho)\succeq qI$. With $\gamma=(1-\eta)/2$, the effect eigenvalues lie in $[\gamma,1-\gamma]$. Thus every report probability lies in this interval, and every normalized output state satisfies
\[
 \rho'\succeq \frac{q\gamma}{1-\gamma}I=aI.
\]
The protected Pauli outcome has variance $1-\tr(\sigma_z\rho')^2\ge4a(1-a)$. Its mean has absolute value at most one. Multiplication by $T$ gives exactly the same calibration constants as in Theorem~\ref{thm:classical}. From the specified preparation, the first $V$ is fair, because $D_q(I/2)=I/2$ and $\tr K_{a,v}^2=1$. Thus the reference report law is physically obtained, not assigned to a formal orbit.

Positive instrument operations on Hermitian matrices, and Born audit effects, give an unnormalized finite linear representation directly. Include a tomographically complete finite set of probe directions among executable future tests; their probabilities separate density operators on the attained image when $q<1/2$. A Borel realization is the compact density-operator body with the visible $Y,T$ audit registers; a redundant direction at a degenerate parameter is factored out. Report total variation is bounded by half the input trace distance after the channel. Normalization has denominator at least $\gamma$, so predictive update and report continuity are explicit on the invariant compact state body. No infinite-dimensional operator closure is used.

If two directions are nonparallel, the associated effects generally fail to commute. For instance, choosing $n_a$ on the $x$ and $z$ axes gives effects with nonzero commutator. Coherence can persist because $q<1/2$ and the effects are invertible; the operation is not a measure-and-prepare reset. The external observer nonetheless has only the charged classical labels. The physical qubit is not available as free observer memory. Finally the acquired $(Y,T)$ law proves \eqref{eq:rectangle-mass}, independently of the classical proof.
\end{proof}

The reliability mark is a measurable experimental quantity, not a promise that the decoder knows the hidden residual. In both models, replacing all past histories by the last visible proxy generally changes the posterior. What survives is the two-sided conditional risk comparison, uniformly after every rejected call. That is the hypothesis verified here and used by the general theorem.

\section{Nonregenerative Gaussian acquisition and a singular stopping transition}
\label{sec:gaussian}
The preceding theorem does not require recovery of the audit mean from the current report. The following consequence makes that distinction quantitative in an informative experiment without an auxiliary reliability coordinate.

Let $X_t\in\{-1,1\}$ be the protected sign at call $t$. Given the surviving history, before the report suppose
\begin{equation}\label{eq:prior-floor}
 q\le \Pp(X_t=1\mid\mathcal H_{t-1})\le1-q,
 \qquad0<q\le\tfrac12.
\end{equation}
Conditionally on $X_t$, report $Y_t\sim N(\mu X_t,1)$, $\mu>0$, and audit $X_\tau$ after stopping. Every call costs one. This hypothesis follows from the raw controlled two-state transition kernel with all entries in $[q,1-q]$. Transitions may depend on the chosen action and on an unknown fixed model. They need not be symmetric. Preparation may be arbitrary because transition precedes reporting. The observer cannot access $X_t$.

Write $f_\pm(y)=\varphi(y\mp\mu)$ and set
\[
 \nu(dy)=\tfrac12(f_++f_-)dy,\qquad
 A(y)=\frac{\min(f_+,f_-)}{f_++f_-}
     =\frac1{1+e^{2\mu|y|}}.
\]
The symmetric fair preparation gives the actual reference law $\nu$. For $t\ge1$, let $h_t\ge0$ be the unique solution of
\begin{equation}\label{eq:two-sided-threshold}
 \overline\Phi(h_t-\mu)+\overline\Phi(h_t+\mu)=t^{-1},
 \qquad Q_\mu(t)=t\overline\Phi(h_t+\mu).
\end{equation}
In particular $h_1=0$. The quantity $Q_\mu(t)$ is the marked profile with proxy $\operatorname{sign}(y)$ and at least two centers: it equals
$t\int_{|y|\ge h_t}A(y)\nu(dy)$. Rearrangement by the increasing reliability $|y|$ proves this identity even for fractional selection rules.

\begin{theorem}[All-history Gaussian acquisition bounds]\label{thm:gaussian-all}
For every $N\ge1$, every checkpoint dictionary size $m\ge2$, and every autonomous budget $M\ge3$,
\begin{equation}\label{eq:gaussian-sandwich}
 4qQ_\mu(2(1-q)N)
 \le R_\cp^\Theta(N,m),\ R_\aut^\Theta(N,M)
 \le4Q_\mu(N).
\end{equation}
The lower bound permits all history-dependent randomized stopping rules and all history-dependent terminal predictions, without a finite-memory constraint. One common three-state program attains the displayed upper, independently of the unknown transitions. Consequently
\begin{equation}\label{eq:gaussian-order}
 R_\cp^\Theta(N,m)\asymp R_\aut^\Theta(N,M)
 \asymp Q_\mu(N)
 \asymp \exp\{-2\mu\sqrt{2\log N}-2\mu^2\}
\end{equation}
as $N\to\infty$, uniformly in the stated $m,M$. The last comparison is not a ratio-one assertion for the risk. The profile itself satisfies
$Q_\mu(N)/\exp\{-2\mu\sqrt{2\log N}-2\mu^2\}\to1$.
\end{theorem}
\begin{proof}
Fix any legal procedure with $\E\tau\le N$, and condition also on its private past. At a surviving history write $p=\Pp(X_t=1\mid\mathcal H_{t-1})$ and let $u_t(y)\in[0,1]$ be its conditional stop probability. Any terminal decoder incurs at least conditional variance. Multiplying that variance by the raw report density gives
\[
 \frac{4p(1-p)f_+(y)f_-(y)}{pf_+(y)+(1-p)f_-(y)}
 \ge4q(1-q)\min(f_+(y),f_-(y)).
\]
Define the actual predictable counting density
$g(y)=\E\sum_{t\le\tau}u_t(y)$. Tonelli gives $0\le g\le N$. Since the conditional report density is at most $2(1-q)$ times $\nu$, terminal mass one implies
$\int g\,d\nu\ge[2(1-q)]^{-1}=:a_0$. Therefore
\[
 R\ge8q(1-q)\int A g\,d\nu
 \ge8q(1-q)a_0 Q_\mu(N/a_0).
\]
The last inequality selects the smallest values of $A$ with a density bounded by $N$ and mass at least $a_0$. This is a lower bound over all actual histories, not a stationary posterior approximation.

For the upper, stop when $|Y_t|\ge h_N$ and predict $\operatorname{sign}(Y_t)$, using one live and two terminal labels. Under either hidden sign the stop probability is the same number $1/N$. Hence it is $1/N$ after every surviving history, and $\E\tau=N$ exactly, although the hidden state does not regenerate. Under either sign the probability of an erroneous accepted sign on a single call is $\overline\Phi(h_N+\mu)$. Squared loss is four on this event. Summing survival probabilities gives risk $4N\overline\Phi(h_N+\mu)$.

To verify the rate, put $x=h_N-\mu$. Mills' inequalities, obtained by integrating $\varphi'(x)=-x\varphi(x)$, give for $x>0$
\[
 \frac{x}{x^2+1}\varphi(x)<\overline\Phi(x)<\frac{\varphi(x)}x.
\]
They imply $x-\sqrt{2\log N}\to0$ and
\[
 \frac{\overline\Phi(x+2\mu)}{\overline\Phi(x)}
 =(1+o(1))\exp(-2\mu x-2\mu^2),\qquad
 N\overline\Phi(x)\to1.
\]
This proves the stated ratio-one asymptotic for $Q_\mu$. For any fixed $c\ge1$, $Q_\mu(cN)/Q_\mu(N)\to1$. Positivity and continuity on bounded intervals then bound this ratio away from zero for all $N\ge1$. Apply this with $c=2(1-q)$. It proves the uniform comparability claims, not an exact finite-budget checkpoint optimum.
\end{proof}

The calibration mechanism is explicit here. If $e$ is the posterior probability that the proxy sign is wrong, then $e\le(1-q)A/q$, while $e(1-e)$ is bounded below by a positive multiple of $A$. Thus $v=4e(1-e)\asymp A$ and $|m-\operatorname{sign}(y)|^2=4e^2\le C A$, uniformly over all histories. The proof above improves the general envelope constants by using the unnormalized variance directly. The protected audit is fixed; selecting a reliable current sign does not estimate an unchanging hidden parameter.

\begin{proposition}[A noncommuting Gaussian instrument]\label{prop:gaussian-quantum}
Let $D_q$ be the qubit channel of Theorem~\ref{thm:quantum}, and for selectable unitaries $V_a$ let
\[
 D_y=\operatorname{diag}(\sqrt{f_+(y)},\sqrt{f_-(y)}),\qquad
 \mathcal G_a(dy)(\rho)=D_yV_aD_q(\rho)V_a^*D_y\,dy.
\]
After stopping, measure $\sigma_z$ without an intervening transition. Then Theorem~\ref{thm:gaussian-all} applies to this audit for every causal selection of $a$. The measurement effects for different $V_a$ can be noncommuting. No common eigenbasis for the instruments or reset of the physical state is assumed.
\end{proposition}
\begin{proof}
The displayed completely positive instrument integrates to a trace-preserving map because $\int D_y^2dy=I$. Before measurement $V_aD_q(\rho)V_a^*\succeq qI$, so its diagonal weights belong to $[q,1-q]$. The joint law of $y$ and the subsequent $\sigma_z$ outcome is precisely $pf_+(y)dy$ or $(1-p)f_-(y)dy$. Thus the raw prior-floor and conditional audit identities used in the proof hold at every history, including histories selected by rejection. Coherences need not vanish after $D_y\rho D_y$, and distinct rotated effects need not commute. Report total variation is Lipschitz in the input trace norm. On bounded report sets the normalized update is uniformly continuous; Gaussian tails remain part of the experiment and are not conditioned away. The finite Hermitian instrument representation gives an executable predictive carrier, and the observer uses only its three charged classical labels.
\end{proof}

\subsection{Actual revealing mass and a singular threshold}
Modify either Gaussian experiment as follows. Independently at each call, with probability $\beta\in[0,1]$ execute a marked revealing operation and report the protected sign exactly. Otherwise execute the Gaussian operation. For the classical model this reports the current sign. For the quantum model it is the rank-one projective instrument $\rho\mapsto P_\pm V_aD_q(\rho)V_a^*P_\pm$, with its sign and branch marked. The immediate audit has that same sign; the next call still starts with its declared transition or channel. This is a genuine mixture of continuous and singular acquired information on the same physical carrier.

\begin{theorem}[A non-power singular acquisition threshold]\label{thm:reveal}
With $m\ge2$ and $M\ge3$, zero risk is possible exactly when $\beta N\ge1$. If $r=1-\beta N>0$ and
\[
 t=\frac{N(1-\beta)}r\ge1,
\]
then
\begin{equation}\label{eq:reveal-sandwich}
 4q\,rQ_\mu(2(1-q)t)
 \le R_\cp^\Theta(N,m),\ R_\aut^\Theta(N,M)
 \le4rQ_\mu(t).
\end{equation}
All constants are uniform in $N,\beta,m,M$ for fixed $q,\mu$. In particular, for fixed $0<\beta<1$ and $N\uparrow1/\beta$, the matching risk order is
\[
 r\exp\{-2\mu\sqrt{2\log(1/r)}\}.
\]
The factor $r$ is actual nonrevealing terminal mass; it is not removed by renormalizing the continuous stratum.
\end{theorem}
\begin{proof}
The expected acquired count of revelations is $\beta\E\tau\le\beta N$. Hence every procedure has nonrevealing terminal mass at least $r$. Apply the preceding unnormalized-variance calculation only on nonrevealing calls. Its counting density is bounded by $N(1-\beta)$, and its $\nu$-mass is at least $r/[2(1-q)]$. Rearrangement gives the lower in \eqref{eq:reveal-sandwich}. This proves strict positivity when $r>0$, including for decoders with unlimited memory.

For the upper, accept every revelation and accept a Gaussian report when $|y|\ge h_t$. The stop hazard at every history is
$\beta+(1-\beta)/t=1/N$. The mean cost is exactly $N$, and the actual nonrevealing terminal mass is $N(1-\beta)/t=r$. Its erroneous-sign contribution is $4N(1-\beta)\overline\Phi(h_t+\mu)=4rQ_\mu(t)$. One live and two terminal labels suffice. If $\beta N\ge1$, wait for the first revelation, of mean cost $1/\beta$, and predict its sign. The asymptotic follows from Theorem~\ref{thm:gaussian-all}; here $t$ is a fixed positive multiple of $1/r$ up to a factor tending to one. No full-rank density lower bound is imposed on the atomic branch.
\end{proof}

\subsection{The complete two-state curve, including the critical atom}
For this subsection only, take the known symmetric classical transition with flip probability $q$, fair preparation, and let $\kappa=(1-2q)/(2q)$. For a one-live, one-terminal program write $a,b$ for its stop probabilities conditional on signs $+1,-1$. Its mean time and terminal mean are
\begin{equation}\label{eq:two-state-occupation}
 T(a,b)=\frac{2+\kappa(a+b)}{a+b+2\kappa ab},\qquad
 m(a,b)=\frac{a-b}{a+b+2\kappa ab}.
\end{equation}
The upper likelihood-ratio boundary consists, in increasing selectivity, of
\begin{align}
 &(a,b)=(1,b),&&1-\beta\le b\le1,\label{eq:roc}\\
 &a=\beta+(1-\beta)\overline\Phi(h-\mu),\quad
 b=(1-\beta)\overline\Phi(h+\mu),&&-\infty<h<\infty,\notag\\
 &(a,b)=(\beta,0),&&\beta>0.\notag
\end{align}
The first segment randomizes on the negative revelation; the middle accepts the positive revelation and the upper Gaussian tail, but no negative revelation. Empty segments at $\beta=0$ or $1$ are omitted.

\begin{theorem}[Exact all-program two-state curve]\label{thm:two-state}
For $\beta=0$, put $N_0=+\infty$; for $\beta>0$, put
$N_0=2/\beta+\kappa$. At $N=1$ the optimal two-state risk is one. For $1<N<N_0$, there is a unique point $(a_N,b_N)$ on \eqref{eq:roc} with $T(a_N,b_N)=N$, and
\[
 R_\aut(N,2)=1-m(a_N,b_N)^2.
\]
For $N\ge N_0$ the optimum is zero. These statements cover all Borel randomized two-state programs. Thus, for $0<\beta\le1$, the exact zero-risk threshold is $2/\beta+\kappa$ for two states, but $1/\beta$ for three or more states. Both thresholds count the live controller.
\end{theorem}
\begin{proof}
Let $v_\pm$ be the actual expected numbers of reports in the two signs, and $d_+=av_+$, $d_-=bv_-$. Summing and subtracting the killed two-state balance gives
$d_++d_-=1$ and $v_+-v_-=-\kappa(d_+-d_-)$. Solving yields \eqref{eq:two-state-occupation}. Every nonzero routing has positive stop probability after a transition, bounded below by $q(a+b)$, and therefore finite occupation. Its optimal single constant prediction is $m$, with risk $1-m^2$. This occupation calculation is inherited from the unmarked Gaussian experiment~\cite{GTF51}; the atomic boundary and its consequences are new here.

Reflect signs if necessary to have $a\ge b$. At fixed $b$, direct differentiation gives
\[
 \partial_a m=\frac{2b(1+\kappa b)}{(a+b+2\kappa ab)^2},\qquad
 \partial_a T=-\frac{2(1+\kappa b)^2}{(a+b+2\kappa ab)^2}.
\]
Thus maximizing $a$ improves both risk and cost, with the zero-$b$ endpoint treated directly. The likelihood ratio is $+\infty$ on the positive revelation, $e^{2\mu y}$ on the Gaussian branch, and zero on the negative revelation. The elementary Neyman--Pearson comparison~\cite{NP} therefore gives exactly \eqref{eq:roc}, including randomization on the final atom. No threshold restriction has been assumed for the competing programs.

Along this nested family the hitting time increases strictly from one to $N_0$: the endpoint has $a=\beta,b=0$ and \eqref{eq:two-state-occupation} gives $2/\beta+\kappa$. For $\beta=0$ the limit is infinite. The ratio $a/b$ increases strictly, while $a$ decreases, so
$m=(a/b-1)/(a/b+1+2\kappa a)$ increases from zero to one. The negative-atom segment obeys the same conclusion directly. This proves existence, uniqueness and budget saturation before $N_0$; after that budget the positive-revelation rule has risk zero. At $N=1$ stopping is compulsory at the first report, and a single terminal value cannot distinguish its sign. The three-state threshold was proved in Theorem~\ref{thm:reveal} over all histories.
\end{proof}

Theorems~\ref{thm:gaussian-all} and \ref{thm:reveal} provide matching all-history checkpoint orders, while Theorem~\ref{thm:two-state} provides an exact finite-budget two-state value. Neither assertion identifies the exact unrestricted checkpoint value at every finite budget, nor the exact intermediate positive-risk curve for every $M\ge3$. These distinctions are essential in comparing a finite-memory machine with a full-history stopping rule.

\subsection*{Finite arithmetic certificate}
\begin{center}
\small
\begin{tabular}{@{}rrlll@{}}
\toprule
$\beta$ & $N$ & all-history lower & three-state upper & two-state optimum \\
\midrule
$0$ & 2 & $[0.021402, 0.021403]$ & $[0.161245, 0.161246]$ & $[0.821792, 0.821793]$ \\
$0$ & 4 & $[0.008837, 0.008838]$ & $[0.057865, 0.057866]$ & $[0.312884, 0.312885]$ \\
$0$ & 8 & $[0.004265, 0.004266]$ & $[0.025744, 0.025745]$ & $[0.084299, 0.084300]$ \\
$0$ & 32 & $[0.001297, 0.001298]$ & $[0.007152, 0.007153]$ & $[0.014194, 0.014195]$ \\
$1/4$ & 2 & $[0.006228, 0.006229]$ & $[0.042805, 0.042806]$ & $[0.790117, 0.790118]$ \\
$1/4$ & 3 & $[0.000952, 0.000953]$ & $[0.005688, 0.005689]$ & $[0.405411, 0.405412]$ \\
$1/4$ & 4 & $[0.000000, 0.000000]$ & $[0.000000, 0.000000]$ & $[0.157475, 0.157476]$ \\
$1/4$ & 8 & $[0.000000, 0.000000]$ & $[0.000000, 0.000000]$ & $[0.001633, 0.001634]$ \\
$1/4$ & 9 & $[0.000000, 0.000000]$ & $[0.000000, 0.000000]$ & $[0.000000, 0.000000]$ \\
$1$ & 1 & $[0.000000, 0.000000]$ & $[0.000000, 0.000000]$ & $[1.000000, 1.000000]$ \\
$1$ & 2 & $[0.000000, 0.000000]$ & $[0.000000, 0.000000]$ & $[0.750000, 0.750000]$ \\
$1$ & 3 & $[0.000000, 0.000000]$ & $[0.000000, 0.000000]$ & $[0.000000, 0.000000]$ \\
\bottomrule
\end{tabular}
\end{center}
Here $q=1/4$ and $\mu=1$. The lower and upper columns enclose the explicitly proved bound functions, not the unknown exact checkpoint optimum. Every upper uses a certified feasible rational threshold or an exact finite revealing-symbol rule. The two-state column encloses the exact optimum proved in Theorem~\ref{thm:two-state}.

The accompanying certificate encloses the normal tails and threshold roots by rational interval arithmetic. It reports the all-history lower, the three-state upper at the same mean budget, and the exact two-state curve, including the revealed-atom transition. These are finite-parameter checks of the formulas; their all-history and continuum validity comes from the preceding proofs.

\section{Composition in one protected task}\label{sec:composition}
We distinguish a certified comparison of experiments from an attained error mechanism. This avoids treating a permitted approximation tolerance as an unavoidable risk lower bound.

\begin{proposition}[Typed causal comparison]\label{prop:composition}
Suppose a parameter-independent causal simulator from a physical experiment $\mathfrak E$ to a target $\mathfrak F$ includes a report transformation, a strategy lift, action compatibility, a stopping handshake and the protected audit readout. Suppose it uses at most $s$ persistent labels and, for every lifted strategy, (i) its complete stopped-and-scored law differs from the target by at most $\varepsilon$ in total variation, and (ii) the expected physical cost is at most $a+b\E\tau$, where $a,b\ge0$ are fixed. Then any target $M$-state observer of expected budget $N$ gives a physical observer with at most $sM$ states, expected cost at most $a+bN$, and risk at most
\[
 R_{\mathfrak F}+L\varepsilon
\]
for a common loss in $[0,L]$. In the same execution law, a terminal decoder perturbation bounded in $L^2$ by $e$ satisfies
\[
 \sqrt{R_{\rm perturbed}}\le\sqrt{R_{\rm exact}}+e.
\]
Independent stage simulation defects may be summed when they are certified conditionally on every surviving history and their expected accumulated sum is bounded. Persistent simulator states multiply; separately stored program bits and workspaces add under a sequential implementation. A reverse risk bound requires a separately certified reverse morphism.
\end{proposition}
\begin{proof}
Run the target label and simulator label as their Cartesian product. The simulator's strategy lift uses no true model parameter, and its terminal handshake precedes the protected audit. The assumed complete-law comparison bounds the expectation of a loss in $[0,L]$ by $L\varepsilon$. The stated cost bound is a separate hypothesis and is not inferred from total variation. Conditional maximal couplings, or telescoping conditional integral differences, bound accumulated bounded-task error by the expected sum of the stage defects. Minkowski's inequality proves the decoder estimate on the same law. A one-way construction says nothing about procedures that do not factor through it, so it cannot prove a reverse minimax comparison.
\end{proof}

Numerical changes to an acceptance rule can change its stopping time. A decoder perturbation alone does not. In particular, closeness of a bounded scored law does not control the mean of an unbounded stopping time. Appendix~\ref{app:compute} gives a separate hazard-margin implementation certificate.

\begin{theorem}[An attained joint resource--error law]\label{thm:joint}
Take either raw model of Section~\ref{sec:realizations}, with $d,\alpha>0$, and augment its protected audit by two independent coordinates. A compulsory preparation call creates a fair sign $U$; a subsequent paid early report is $U$ with probability $1-2\varepsilon$ and the erasure symbol otherwise, where $0\le\varepsilon\le1/2$. A fixed calibration alternative $\zeta\in\{-\delta,\delta\}$ is unobservable in every report. The core experiment and its model are independent of $(U,\zeta)$; the terminal audit is $(Z,U,\zeta)$. The prediction is committed before this audit. Squared losses of the three coordinates are added.

Count preparation, early report and final audit as three physical calls, and put $N=n+3$ with $n\ge1$. Then, for every $M\ge7$ and bounded $\delta$,
\begin{equation}\label{eq:joint-law}
 \mathcal R_{\rm aut}(N,M,\delta,\varepsilon)
 \asymp F_{d,\alpha}(n,M-1)+\delta^2+\varepsilon.
\end{equation}
The constants are independent of $n,M,\delta,\varepsilon$ on the stated range. The early erased-bit experiment has causal deficiency exactly $\varepsilon$ relative to the perfect early-bit interface, whose virtual report must be completed before any core action. Thus the last term is attained, not merely an allowed error ceiling.
\end{theorem}
\begin{proof}
For the lower bound, permit the core controller the entire history, including the early outcome and every fresh coin. The calibration certificate still holds after each of these histories, so the actual core occupation proof with its total expected budget $n$ remains valid. At most $M-1$ terminal core predictions can be used. This gives $cF_{d,\alpha}(n,M-1)$, without assuming that the three early-report lanes receive equal conditional budgets. On the erasure event $U$ remains independent of all subsequently observed core data and routing coins; its Bayes squared loss is one there. Its contribution is at least $2\varepsilon$. Averaging the two calibration alternatives gives at least $\delta^2$, since all report laws are identical under them. These three conditional lower bounds add for the same execution, and averaging is a valid lower bound to the worst calibration alternative and model.

For the upper, put $m=\lfloor(M-4)/3\rfloor\ge1$. Use one initial early-report label, three copies of the core live label, and $3m$ terminal labels. The total is $4+3m\le M$. Preparation occurs before this initial observer label and is nevertheless charged physically; the audit occurs only after terminal commitment and is also charged. Each lane uses the same core program of Theorem~\ref{thm:geometry}. Read out the known bit on the two revealed lanes and zero on the erasure lane; always read out zero for calibration. The risk is at most $CF_{d,\alpha}(n,m)+2\varepsilon+\delta^2$. Since $m$ and $M-1$ differ by a uniformly bounded factor for $M\ge7$, this is \eqref{eq:joint-law}.

For deficiency, regard the two possible bit preparations $U=\pm1$ as the experiment parameter. A causal simulator copies an available bit and uses a fresh fair sign on erasure. Its worst parameter total variation from the perfect early bit is $\varepsilon$. Conversely, under the fair prior, no transformation of an erased bit can guess its sign with error less than $\varepsilon$. The indicator test for an incorrect virtual early report is bounded by total variation, proving the opposite inequality. The prescribed early completion prevents using the final audit to manufacture that report retrospectively. Later core execution can be coupled on the correct virtual-bit event. This proves exact deficiency for the declared complete causal interface.
\end{proof}

The same construction attaches the real $\delta^2$ and $\varepsilon$ terms to the Gaussian and revealing-mass bounds, with the same lane accounting; the term for the core is then the corresponding all-history lower and explicit three-state upper rather than an asserted exact checkpoint value. An additive numerical error requires the independent execution certificate just stated. Theorems~\ref{thm:central}, \ref{thm:geometry}, \ref{thm:reveal} and \ref{thm:joint} therefore distinguish geometric distortion, acquisition uncertainty, realized calibration ambiguity, numerical implementation and causal comparison.

\appendix
\section{Predictive quotient and experiment-level meaning}\label{app:quotient}
We record the measurable foundation needed by the raw realizations, without identifying a selected task proxy with the whole predictive state. Let $\mathcal X$ be a standard Borel prepared predictive carrier and let $(T_i)_{i\ge1}$ be a countable determining family of executable finite future tests. It includes report-cylinder tests and the permitted task readouts, and is closed under the declared causal interfaces. ``Determining'' means that equality of these test probabilities is equivalent to equality for the specified legal future-test class. Write
\[
 e(x)=(\Pp_x(T_i=1))_{i\ge1}\in[0,1]^{\mathbb N}.
\]
Report kernels and conditional updates are assumed Borel on the carrier and behaviorally compatible: if $e(x)=e(x')$, their next report laws agree and their updated $e$-coordinates agree almost everywhere under that common report law. These properties can be verified by the raw unnormalized instruments rather than inferred from a task covariance.

\begin{proposition}[Measurable minimal realization]\label{prop:quotient}
Suppose $e$ is Borel, its image $Q$ is Borel, and it has a Borel section $s:Q\to\mathcal X$. Then $Q$, with the report and update kernels descended through $s$, is a standard Borel realization of future-test equivalence. It is minimal in the following precise sense: for any other standard Borel statistic $T:\mathcal X\to Z$ with Borel image and section, if $T(x)=T(x')$ implies equality of all future tests, there is a unique Borel map $g:T(\mathcal X)\to Q$ with $e=g\circ T$. An actual acquisition law is pushed forward by $e$; its support is not replaced by all of $Q$.

The image and section hypotheses hold for a compact metrizable carrier with continuous test probabilities. They are not automatic for an arbitrary Borel image.
\end{proposition}
\begin{proof}
The coordinates of $e$ show that its fibers are exactly predictive equivalence classes. Compose the original report law and the $e$-valued update with $s$. Behavioral compatibility makes the descended law independent of the representative, and composition makes it Borel. If $s_T$ is a section of $T$, set $g=e\circ s_T$; it is Borel and is constant on each fiber by the assumed sufficiency of $T$. Surjectivity onto $T(\mathcal X)$ gives uniqueness.

In the compact continuous case $Q$ is compact. To construct a Borel section, take successive finite covers of the compact carrier by closed sets of diameter tending to zero. At each stage choose the first set whose intersection with the previously selected sets and the fiber is nonempty. Nonemptiness is a closed condition on $Q$, since the relevant compact set has compact continuous image. The choices are Borel, nested compact intersections are nonempty and have vanishing diameter, and their unique limit is a Borel selected preimage. This proves the section assertion. It does not assert that every image of a measurable history map is Borel.
\end{proof}

For an unknown fixed model family, the theoretical carrier retains the model coordinate: it is $(\theta,x)$, not $x$ alone. Future-test evaluation uses the same fixed $\theta$ throughout. This coordinate is not supplied to the observer or to its common program. A compact model family with jointly continuous test probabilities again satisfies the compact selection argument; a merely Borel family must supply the stated image and section hypotheses separately. Equality of current physical density operators across different models is not, by itself, future-test equivalence.

Finite positive classical instruments and finite-dimensional positive quantum instruments supply the unnormalized carriers in Section~\ref{sec:realizations}. Distinct physical coordinates are retained only when executable tests separate them. Redundant directions, including degeneracies, are factored out. Continuous finite probe families on the compact belief or density-operator body give the stated realization; visible audit registers are included when needed. In the Gaussian examples, current-report continuity is local on bounded report sets, and tail probabilities are kept in the actual experiment. No unbounded-operator assertion is involved.

When equivalence has been verified for all finite paths under a given legal strategy class, bounded stopped tests with almost-surely finite stopping are obtained by truncation: the difference between a stopped test and its finite truncation is bounded by a constant times $\Pp(\tau>n)$, which tends to zero. This is an observational equivalence assertion, not a finite-memory implementation of the truncation clock or a uniform infinite-horizon resource estimate. Every theorem in the main text proves its own stopped occupation and cost bound.

The marked acquired law is the pushforward of an actual reference report law by a task-specific proxy and uncertainty mark. It is an instrument for comparing risks under the raw calibration certificate. It need not parameterize $Q$, and no likelihood-curvature, Fisher, covariance, or dimensional formula is used to manufacture predictive minimality.

\section{Finite-description execution and separate resource charges}\label{app:compute}
The Borel program theorem and its finite-bit implementations are distinct assertions. Arbitrary Borel kernels or noncomputable measures do not by themselves provide executable numerical primitives. For a finite implementation assume supplied, parameter-uniform certificates for: rational descriptions of the selected codebook and action probabilities; evaluation of the raw reports and accept regions with controlled probability of an incorrect decision; sampling coins to prescribed variation accuracy; and terminal readout error. A packing or cover used by the geometric theorem must have a computable finite certificate. Products of intervals and Cantor cylinders have such explicit lists. No minimization oracle is presumed to run for free.

\begin{proposition}[Hazard-margin finite execution]\label{prop:finite}
Consider a one-live-label program whose conditional stop probability at every history is at least $p>0$, whose per-call contribution to a loss in $[0,L]$ is at most $b$, and whose mean cost is controlled by its fixed per-call action-cost mean $c$. Suppose a finite-description implementation changes the conditional joint action, retained decision and audit-relevant state law by at most $\eta<p$ in total variation, with the raw report coupled identically, uniformly in every history and model. Then
\[
 \E\tau_{\rm num}\le\frac1{p-\eta},\qquad
 R_{\rm num}\le\frac{b+L\eta}{p-\eta},\qquad
 \E C_{\rm num}\le\frac{c+c_{\max}\eta}{p-\eta}.
\]
Terminal decoder root error adds as in Proposition~\ref{prop:composition}. The same label alphabet suffices. Under the above computational certificates, the geometric matching upper can be implemented by finite code and finite current-call precision without changing its $N,m$ order, while meeting the original mean budget exactly as an inequality.
\end{proposition}
\begin{proof}
This is a decision/routing coupling; it does not require total-variation convergence of a discrete numerical report to the continuous raw report. The implemented stop hazard is at least $p-\eta$, so its actual survival sum is at most $(p-\eta)^{-1}$. Its conditional stopped-loss contribution is at most $b+L\eta$. Summing these two bounds proves the first two assertions; the same argument for expected per-call cost proves the third. These are direct hazard estimates, not expected-time conclusions from a terminal total-variation estimate.

For unit-cost geometric programs at large $N$, build the reference upper with budget $N/2$. Its certified hazard is at least $2/N$. Choose the one-call decision error at most $1/N$ and, more stringently, at most a constant multiple of $F_{d,\alpha}(N,m)/N$. The displayed inequalities then give mean at most $N$ and risk of order $F_{d,\alpha}(N,m)$, since this profile is stable under fixed budget rescaling. Choose readout error of order $\sqrt{F_{d,\alpha}(N,m)}$. Bounded remaining budgets use the fixed one-call decoder. Costs other than one are treated by a strict cost margin in the displayed third inequality. At singular thresholds where rescaling changes the phase, this argument is not used: a specific certified feasible threshold or the exact revealed-atom policy is required instead.
\end{proof}

A program with $m$ centers of coordinate precision $p_r$, finitely many accept-region descriptions of total length $L_A$, and rational action/coin data of length $L_P$ has description length bounded by
$O(kmp_r+L_A+L_P+\log m)$ for a fixed instruction set. Persistent labels require $\lceil\log_2(m+1)\rceil$ bits, not the entire preceding description length. Current-call evaluation requires separately declared workspace $W(p)$ and time $T(p)$. At fixed implemented precision the expected total evaluation time is at most $\E\tau_{\rm num}\,T(p)$, plus preparation, planning and audit costs. A finite expected call count gives no deterministic running-time bound.

For Gaussian tails, overflow regions and threshold-boundary probabilities must both be included in the uniform decision certificate; clipping away their mass is not permitted. The rational tail certificates accompanying Section~\ref{sec:gaussian} certify feasible thresholds and risks, not a cost-free infinite-precision normal report. At a reveal-only zero-risk point, reading the finite report symbol is exact and its finite description is immediate. For a continuous threshold, rational feasible approximants approach the stated value; an exact mathematical threshold is not described as a finite binary real. Program bits, temporary arithmetic space, calibration precision, actual calls, wall-clock primitive time and simulator state therefore remain distinct entries in the resource ledger.

\begin{thebibliography}{99}
\bibitem{GTF32} Q. Qi, \emph{General Theta Foundations I: Acquired Geometry and Causal Resource Transfer}, Bayesian-curvature article, 2026. Complete technical companion.
\bibitem{GTF49} Q. Qi, \emph{General Theta Foundations I: Acquired Geometry and Causal Resource Transfer}, acquired-risk article, 2026. Complete technical companion.
\bibitem{GTF51} Q. Qi, \emph{General Theta Foundations I: Acquired Geometry and Causal Resource Transfer}, causal-acquisition article, 2026. Complete technical companion.
\bibitem{Franc} V. Franc, D. Prusa and V. Voracek, Optimal strategies for reject option classifiers, \emph{J. Mach. Learn. Res.} \textbf{24} (2023), no.~11, 1--49.
\bibitem{CGM} J. A. Cuesta-Albertos, A. Gordaliza and C. Matr\'an, Trimmed $k$-means: an attempt to robustify quantizers, \emph{Ann. Statist.} \textbf{25} (1997), 553--576. doi:10.1214/aos/1031833664.
\bibitem{HC} M. E. Hellman and T. M. Cover, Learning with finite memory, \emph{Ann. Math. Statist.} \textbf{41} (1970), 765--782. doi:10.1214/aoms/1177696958.
\bibitem{SC} C. R. Shalizi and J. P. Crutchfield, Computational mechanics: pattern and prediction, structure and simplicity, \emph{J. Statist. Phys.} \textbf{104} (2001), 817--879.
\bibitem{PB} P. Poupart and C. Boutilier, Value-directed compression of POMDPs, \emph{Adv. Neural Inf. Process. Syst.} \textbf{15} (2002), 1547--1554.
\bibitem{Varshney} K. R. Varshney, Bayes risk error is a Bregman divergence, \emph{IEEE Trans. Signal Process.} \textbf{59} (2011), 4470--4472. doi:10.1109/TSP.2011.2159500.
\bibitem{Kim} M. Kim, \emph{Finite Memory Belief Approximation for Optimal Control in Partially Observable Markov Decision Processes}, arXiv:2601.03132v1, 2026.
\bibitem{Rabiner} L. R. Rabiner, A tutorial on hidden Markov models and selected applications in speech recognition, \emph{Proc. IEEE} \textbf{77} (1989), 257--286.
\bibitem{Jacobs} K. Jacobs and D. A. Steck, A straightforward introduction to continuous quantum measurement, \emph{Contemp. Phys.} \textbf{47} (2006), 279--303. doi:10.1080/00107510601101934.
\bibitem{NP} J. Neyman and E. S. Pearson, On the problem of the most efficient tests of statistical hypotheses, \emph{Philos. Trans. Roy. Soc. London Ser. A} \textbf{231} (1933), 289--337. doi:10.1098/rsta.1933.0009.
\end{thebibliography}

\end{document}
